\documentclass[11pt]{article}

%%%%%%%%%%%%%%%%%%%%%%%%%%%%%%%%%%%%%%%%%%%%%%%%%%%
% Preamble:
%%%%%%%%%%%%%%%%%%%%%%%%%%%%%%%%%%%%%%%%%%%%%%%%%%%
% Typical Packages:
\usepackage{fullpage}
\usepackage{amsfonts}
\usepackage{amsmath}
\usepackage{amsthm}
\usepackage{graphicx}
\usepackage{color}
\usepackage{amssymb}
\usepackage{empheq}
\usepackage{mathrsfs}
\usepackage{enumerate}
\usepackage{tikz}
\usepackage{upgreek}
\usepackage{tipa}
\usepackage{multicol}
\usepackage{verbatim}
\usepackage{floatrow}
\usepackage{gensymb}
\usepackage{caption}
\usepackage[T1]{fontenc}
\usepackage[font=small,labelfont=bf,tableposition=top]{caption}

\thispagestyle{empty}

%%%%%%%%%%%%%%%%%%%%%%%%%%%%%%%%%%%%%%%%%%%%%%%%%%%%
% Start Document:
%%%%%%%%%%%%%%%%%%%%%%%%%%%%%%%%%%%%%%%%%%%%%%%%%%%%
\begin{document}

% Title:
\begin{center}
\Large{\textsc{Pre-Course Survey }}\\
\end{center}

%\vspace{.2cm}


%%%%%%%%%%%%%%%%%%%%%%%%%%%%%%%%%%%%%%%%%%%%%%%%%%%%
% Demographics, questions 1-4
%%%%%%%%%%%%%%%%%%%%%%%%%%%%%%%%%%%%%%%%%%%%%%%%%%%%
\begin{enumerate}[{$\qquad 1.]$}]
\item What is your name?$\line(1,0){38}\line(1,0){195}$
\item What is your gender?$\line(1,0){38}\line(1,0){195}$
\item What is your major?$\line(1,0){38}\line(1,0){195}$
\item Circle your class standing: freshman, sophomore, junior, or senior. 


%%%%%%%%%%%%%%%%%%%%%%%%%%%%%%%%%%%%%%%%%%%%%%%%%%%%
%The following questions (5-7) will asses SLO 2: Quantitative Reasoning in Calculus
%%%%%%%%%%%%%%%%%%%%%%%%%%%%%%%%%%%%%%%%%%%%%%%%%%%%
\item Consider the following statement: ``I feel confident applying mathematical concepts from calculus to solve quantitative problems.''
\begin{center} 
(1) Strongly Agree, \qquad (2) Agree, \qquad (3) Disagree, \qquad (4) Strongly Disagree. \end{center}

\item Consider the following statement: ``I can approach a calculus problem in more than one way (for example, algebraically, graphically, or numerically).''
\begin{center} 
(1) Strongly Agree, \qquad (2) Agree, \qquad (3) Disagree, \qquad (4) Strongly Disagree. \end{center}

\item Consider the following statement: ``When faced with a challenging calculus problem, I am confident I can find a way to solve it.''
\begin{center} 
(1) Strongly Agree, \qquad (2) Agree, \qquad (3) Disagree, \qquad (4) Strongly Disagree. \end{center}


%%%%%%%%%%%%%%%%%%%%%%%%%%%%%%%%%%%%%%%%%%%%%%%%%%%%
%The following questions (8-13) will asses SLO 5: Familiarity, Proficiency, and Feelings towards Techonology
%%%%%%%%%%%%%%%%%%%%%%%%%%%%%%%%%%%%%%%%%%%%%%%%%%%%
\item Consider the following statement: ``In past math courses, using technology helped me to stay engaged with the material.''
\begin{center}
(1) Strongly Agree, \qquad (2) Agree, \qquad (3) Disagree, \qquad (4) Strongly Disagree. \end{center}

\item Consider the following statement: ``I believe technology can make learning calculus less intimidating.''
\begin{center} 
(1) Strongly Agree, \qquad (2) Agree, \qquad (3) Disagree, \qquad (4) Strongly Disagree. \end{center}

\item Consider the following statement: ``I believe technology can help me better understand abstract calculus concepts.''
\begin{center} 
(1) Strongly Agree, \qquad (2) Agree, \qquad (3) Disagree, \qquad (4) Strongly Disagree. \end{center}

\item Consider the following statement: ``I believe technology can be used to visualize and experiment with mathematical ideas.''
\begin{center} 
(1) Strongly Agree, \qquad (2) Agree, \qquad (3) Disagree, \qquad (4) Strongly Disagree. \end{center}

	\pagebreak 

\thispagestyle{empty}

\item Consider the following statement: ``I am confident in my ability to learn and use new technological tools to help me understand mathematical ideas.''
\begin{center} 
(1) Strongly Agree, \qquad (2) Agree, \qquad (3) Disagree, \qquad (4) Strongly Disagree. \end{center}

\item Consider the following statement: ``I am confident using graphing technology (e.g.~graphing calculators, online tools) to explore mathematical problems.''
\begin{center} 
(1) Strongly Agree, \qquad (2) Agree, \qquad (3) Disagree, \qquad (4) Strongly Disagree. \end{center}


%%%%%%%%%%%%%%%%%%%%%%%%%%%%%%%%%%%%%%%%%%%%%%%%%%%%
%The following questions (14-17) will asses Desmos Proficiency and Familiarity:
%%%%%%%%%%%%%%%%%%%%%%%%%%%%%%%%%%%%%%%%%%%%%%%%%%%%
\item Have you ever used Desmos yourself for homework, projects, or exploring math on your own? 
\begin{center} 
(1) Yes, frequently, \qquad (2) Yes, occasionally, \qquad (3) Rarely, \qquad (4) Never. \end{center}

\item Consider the following statement: ``I am confident using Desmos to perform more advanced tasks (for example, animations, tables, etc) other than just graphing functions or calculating numerical values.''
\begin{center} 
(1) Strongly Agree, \qquad (2) Agree, \qquad (3) Disagree, \qquad (4) Strongly Disagree. \end{center}

\item Have you ever had a teacher or professor use Desmos in a class to demonstrate mathematical ideas? 
\begin{center} 
(1) Yes, in college, \qquad (2) Yes, in high school, \qquad (3) Yes, in both, \qquad (4) No. \end{center}

\item Have you ever learned from Desmos (for example, by watching demonstrations, tutorials, or visualizations) without actively using the tool yourself? 
\begin{center} 
(1) Yes, \qquad (2) No. \end{center}


%%%%%%%%%%%%%%%%%%%%%%%%%%%%%%%%%%%%%%%%%%%%%%%%%%%%
%The following questions (18-19) are open ended.
%%%%%%%%%%%%%%%%%%%%%%%%%%%%%%%%%%%%%%%%%%%%%%%%%%%%
\item What strategies or tools have been most helpful for you in learning mathematics in the past? \vfill

\item What do you expect (or hope) technology, such as Desmos, will add to your learning experience in this course? \vfill


\end{enumerate}

\end{document}